% \documentclass[addpoints]{exam}
\documentclass[addpoints,12pt]{exam}

\usepackage[slovene]{babel}
\usepackage{amsfonts}
\usepackage{amssymb}
\usepackage{amsmath}
\usepackage{float}
\usepackage{graphicx}
\usepackage{enumitem}
\usepackage{color}
\usepackage[gen]{eurosym}
\usepackage{newunicodechar}
\newunicodechar{€}{\euro}
\usepackage{arev}

%Komentar naslednje vrstice glede na verzijo testa -- rešitve ali prostor za odgovore:
% \printanswers

\linespread{2}


\pointsinrightmargin
\marginbonuspointname{}


\renewcommand{\solutiontitle}{\noindent\textbf{Rešitve in točkovanje:}\par\noindent}

\colorgrids
\definecolor{GridColor}{gray}{0.7}
\setlength{\gridsize}{7mm}

\extrawidth{5mm}
\setlength{\rightpointsmargin}{1.5cm}

\begin{document}

\runningfooter{}{}{}

\begin{center}
    \fbox{\fbox{\parbox{15cm}{\centering
    Popravljanje in pridobivanje ocen -- četrto pisno ocenjevanje znanja \\ 1. letnik -- test G \\ 12. junij 2025 \\ (Realna števila)}}}
\end{center}

\vspace{0.1in}
\makebox[\textwidth]{Ime in priimek, oddelek:\enspace\hrulefill}


\begin{center}
    \hqword{Naloga}
    \hpword{Možne točke}
    \chbpword{Dodatne točke}
    \hsword{Dosežene točke}
    \htword{Skupaj}  
    % \combinedpointtable[h][questions]
    \gradetable[h][questions]

\end{center}

\vspace{0.1in}


\begin{questions}

    % 1. vprašanje

    \question[5]
    Natančno izračunajte vrednost izraza.
    $$ \sqrt{\left(3-2\sqrt{7}\right)^2}+\sqrt[3]{343}-\left(\sqrt{7}+1\right)^2 $$

    \begin{solution}[7cm]
        \begin{align*} 
            \dfrac{\sqrt{5}}{\sqrt{5}+1}-\dfrac{5-3\sqrt{5}}{4\sqrt{5}}+\left(1+\sqrt{5}\right)^2-\sqrt{125}&=\dfrac{\sqrt{5}(\sqrt{5}-1)}{(\sqrt{5}+1)(\sqrt{5}-1)}-\dfrac{(5-3\sqrt{5})\sqrt{5}}{4\sqrt{5}\sqrt{5}}+1+2\sqrt{5}+5-5\sqrt{5}= \\
                                                                                                            &=\dfrac{5-\sqrt{5}}{5-1}-\dfrac{5\sqrt{5}-3\cdot 5}{4\cdot 5}+6-3\sqrt{5}= \\
                                                                                                            &=\frac{5-\sqrt{5}}{4}-\dfrac{\sqrt{5}-3}{4}+6-3\sqrt{5}= \\
                                                                                                            &=\frac{8-2\sqrt{5}}{4}+6-3\sqrt{5}= \\
                                                                                                            &=8-\frac{7}{2}\sqrt{5}= 
        \end{align*}
        Za racionalizacijo $2$ točki, za izračun kvadrata dvočlenika $1$ točka, delno korenjenje $1$ točka, izračun $1$ točka, končne rezultat $1$ točka.
    \end{solution}
    
    
    \newpage

    % 2. vprašanje
    \question[4]
    Rešite enačbo.
    $$ \frac{2}{3}x^2+\frac{4}{3}x=\frac{1}{3}x^2+3x-2 $$
    
    
    \begin{solution}[4cm]
        \begin{align*} 
            \dfrac{2}{x-3}-\dfrac{x}{x-1} &=\dfrac{2x}{x^2-4x+3} \\
            2(x-1)-x(x-3) &= 2x \\
            2x-2-x^2+3x-2x &=0 \\
            x^2-3x+2 &=0 \\
            (x-1)(x-2) &=0
        \end{align*}
        Pravilno odpravljanje ulomkov oziroma skupni imenovalec $1$ točka, zapis in upoštevanje pogojev za rešitev ($x\notin\{1,3\}$) $1$ točka, pravilno reševanje razcepne enačbe $1$ točka, zapis rešitve $x=2$ $1$ točka.
    \end{solution}

    \newpage

    % 3. vprašanje
    \question[4]
    Rešite neenačbo, rešitev grafično upodobite in zapišite z intervalom.
    $$ 4\left(3-\left(2+7x\right)\right)\geq 1-3\left(4x+2\right) $$
    
    
    \begin{solution}[10cm]
        Rešitev prve enačbe $x>3$ ali $x\in(3,\infty)$ $1$ točka, rešitev druge enačbe $x\geq 6$ ali $x\in[6,\infty)$ $1$ točka, rešitev sistema $x\geq 6$ ali $x\in[6,\infty)$ $1$ točka, zapis končne rešitve z intervalom $1$ točka, grafična upodobitev $1$ točka.
    \end{solution}

\newpage
    % 4. vprašanje
    \question[4]
    Rešite sistem enačb.
    $$ \begin{aligned}
        -4x-3y+1&=0 \\ 3x+2y&=-5
    \end{aligned} $$
    
    
    \begin{solution}[8cm]
        Pravilno reševanje sistema enačb do $2$ točki, pravilno sklepanje o neobstoju rešitve in zapis tega do $2$ točki.
    \end{solution}


\newpage

    % 5. vprašanje
    \question[4]
    V litrski bučki imamo $400~mL$ $75~\%$ raztopine natrijevega klorida. To bi s prilivanjem čiste vode želeli razredčiti do $24~\%$ raztopine.
    Koliko vode moramo dodati? Ali imamo bučko zadostne prostornine?
    
    \begin{solution}[9cm]
        ???
    \end{solution}




    % \newpage 

    
    % % 6. vprašanje
    % \question[4]
    % Natančno izračunajte.
    % $$ \left\lvert \lvert\sqrt{3}+4\rvert-\sqrt{3}\right\rvert + \left\lvert \lvert\sqrt{3}-4\rvert-\sqrt{3}\right\rvert $$
    
    
    
    % \begin{solution}[10cm]
    %     \begin{align*}
    %         \left\lvert \lvert\sqrt{3}+4\rvert-\sqrt{3}\right\rvert + \left\lvert \lvert\sqrt{3}-4\rvert-\sqrt{3}\right\rvert &=
    %         \left\lvert \sqrt{3}+4-\sqrt{3}\right\rvert + \left\lvert -\sqrt{3}+4-\sqrt{3}\right\rvert= \\
    %         &= \left\lvert 4\right\rvert + \left\lvert 4-2\sqrt{3}\right\rvert= \\
    %         &= 4+4-2\sqrt{3}= \\
    %         &= 8-2\sqrt{3}
    %     \end{align*}
    %     Pravilno razreševanje absolutnih vrednosti do $3$ točke, končen rezultat $1$ točka.
    % \end{solution}



    % 7. vprašanje
    % \question[5]
    % Rešite enačbo in opišite njen geometrijski pomen ter rešitev grafično prikažite.
    % $$  7=\left\lvert x-5\right\rvert $$
    
    
    % \begin{solution}[8cm]
    %     Opis geometrijskega pomena (oddaljenost števila $x$ od števila $-5$ je $4$) $1$ točka, grafična predstavitev $1$ točka, reševanje enačbe z upoštevanjem pogojev do $2$ točke, zapis rešitve $x\in\{-9,-1\}$ $1$ točka..
    % \end{solution}


    % \newpage

    % \makebox[\textwidth]{Ime in priimek, oddelek:\enspace\hrulefill}

    % ~\\


\newpage
    % 8. vprašanje
    \question[7]
    Rešite enačbo in rešitev grafično prikažite.
    $$\left\lvert x+6 \right\rvert + \left\lvert x-4 \right\rvert=12 $$
    
    
    \begin{solution}[10cm]
        ???
    \end{solution}

    \newpage

    % % 9. vprašanje
    % \question[4]
    % V posodi imamo $440~g$ $95~\%$ raztopine kisline. Koliko gramov $60~\%$ raztopine kisline moramo dodati, da bomo dobili $80~\%$ raztopino kisline?
    
    
    % \begin{solution}[10cm]
    %     Zapis enačbe po besedilu $1$ točka, pravilno reševanje enačbe $1$ točka, zapis pravilne rešitve $1$ točka, zapis pravilnega odgovora $1$ točka.
    % \end{solution}




    % 10. vprašanje
    \question[5]
    Osnovno šolo Dragotina Ketteja obiskuje $450$ učencev. Šola se trenutno prenavlja in vodstvo se je odločilo, da bo na vseh hodnikih stala zidna poslikava učencev.
    Učenci morajo poslikati stene dolžine $600~m$. Vsak od $250$ učencev razredne stopnje bo poslikal $1~m$ stene v $6~urah$. Učencev predmetne stopnje je $200$ in vsak poslika svoj del stene v $4~urah$.
    Kolikšno dolžino stene bi poslikal učenec predmetne stopnje v $6~urah$?
    
    
    \begin{solution}[3cm]
        Pravilen postopek reševanja naloge do $2$ točki, zapis pravilne rešitve $1$ točka, zapis pravilnega odgovora $1$ točka.
    \end{solution}

     \newpage



    % 3. vprašanje
    % \question[4]
    % Produkt dveh števil je $192$, njun največji skupni delitelj pa $4$. Poiščite števili.
    
    % \begin{solution}[7cm]
    %     \begin{align*}
    %         1-(1+x^8)(1+x^4)(1+x^2)(1+x)(1-x)&=1-(1+x^8)(1+x^4)(1+x^2)(1-x^2)= \\
    %             &=1-(1+x^8)(1+x^4)(1-x^4)= \\
    %             &=1-(1+x^8)(1-x^8)= \\
    %             &=1-(1-x^{16})= \\
    %             &=1-1+x^{16}= \\
    %             &=x^{16}
    %     \end{align*}
    %     Za pravilno uporabo pravila razlike kvadratov $1$ točka, za pravilno upoštevan negativen predznak $1$ točka, končen rezultat $1$ točka.
    % \end{solution}


    % \begin{description}
    %     \item[a] $(x-a)^2+(y-b)^2-c^2=0$ 
    %     \item[b] $d^2(x-a)^2+f^2(y-b)^2-(df)^2=0$
    %     \item[c] $S(a,b)$
    %     \item[d] $e^2=d^2-f^2;\ d>f$
    % \end{description}
    % Krožnica: \fillin[a c][2cm]
    % \\ Elipsa:  \fillin[b c d][2cm]




    %  \newpage

    %  \noindent Ime in priimek: \underline{~~~~~~~~~~~~~~~~~~~~~~~~~~~~~~~~~~~~~~~~~~~~~~~~}

    %  ~\\

    %dodatno vprašanje
    % \bonusquestion[2]
    % \textit{Dodatna naloga:} Obravnavajte enačbo. $$a(x-a)=4(x-1)-12$$

    % \begin{solution}[10cm]
    %     Pravilna obravnava enačbe do $2$ točki.
    % \end{solution}




    % \newpage
        

    % dodatno vprašanje
    % \bonusquestion[3]
    % \textit{Dodatna naloga:} Natančno izračunajte. $$ \sqrt{\sqrt{\left(\sqrt{2}-1\right)^2}+\sqrt{\left(\sqrt{2}-2\right)^2}} $$

    % \begin{solution}[7cm]
    %     Pravilen izračun do $3$ točke.
    % \end{solution}




\end{questions}



\end{document}